\section{Blood on the Streets}
\label{sec:bloodyStreets}

To reach the monastery, the characters will need to travel through streets of Kholinar. That will spawn multiple opportunities for smaller encounters on the streets. Use this section for some inspiration, especially if the characters take more time to reach the monastery.

However, before you allow the PCs to chart their course, read the following:

\quoteBox{
	When you step on the street, you see a smoke rising from the city walls. Some people run away from that direction. \\
	"The gate! They are attacking the gate!", one of them shouts.
}

From here, the characters can rush towards the city gates, sneak to the monastery or even attempt their luck to enter the palace. When they walk on the streets it is possible for the events below to happen:



\subsection{More Proclamations}
\label{ssec:bloodyStreetsProclamations}

If long enough time has passed since the previous queen’s proclamation read by her herald, a new one can be announced. Each proclamation starts with the same preamble:

\quoteBox{
	We, \npcQueenName{}, by the grace of Almighty and His Heralds, Queen of the Alethkar, do hereby proclaim that...
}

Each new proclamation errs on the edge of insanity, painting the picture of a queen descending into madness and religious fanaticism. None of the announcements seem to acknowledge ongoing riots. Examples may include:

\begin{itemize}
	\item All lights should be put out after dark and the spheres covered with cloth to make it harder for Voidbringers to find Kholinar in the dead of the night.
	
	\item People should pray to the Almighty out in the open each day, including highstorm days, to beg for absolution and mercy for Kholinar.
	
	\item Each household should sacrifice a chull to beg the Almighty for the King's victory over the Voidbringers on the Shattered Plains.
	
	\item All people, both men and women, should adjust their clothing to cover both freehand and safehand as a sign of modesty and humility.
	
	\item All people should face the Origin and sing ten songs to the ten heralds ten times a day.
	
	\item Tallew gathered on odd days of the week is proclaimed impure and should be burnt. Only tallew gathered on even days are admitted for consumption.
	
	\item People shall bring all of their foods and drinks to the ardents to be purified with prayers before they can be consumed.	
\end{itemize}



\subsection{Violent Clash}

When starting this encounter, read the following:

\quoteBox{
	You hear battle roars and cries of pain. The street is red with blood. You can't tell who started this fight, there are wounded on both sides.
}

Choose 2 groups of adversaries:
\begin{itemize}
	\item 1d4 \enemyWallGuard{}s
	\item 1d6 \enemyRioter{}s
	\item 1d6 \enemyCultist{}s (each \enemyCultistOfMoments{} of random kind)
\end{itemize}

Those groups fight with each other. They ignore the characters, but will attack them once they enter the conflict helping one of the sides.


\subsection{Helpless Parshmen}

When starting this encounter, read the following:

\quoteBox{
	You see a docile parshmen worker, hauling some boxes. Angry mob gathers around him. \\
	"Storming Voidbringer!", one of the rioters shouts as he raises his hammer. "Die, you filth!"
}

A \enemyPrshmen{} is surrounded by 1d6 \enemyRioter{}s. Unless the characters intervene, they will kill the worker. The party can also join them and execute queen's order.

You can repeat this encounter multiple times, changing some details of the scene:
\begin{itemize}
	\item You can change aggressors' fraction, cycling between \enemyRioter{}s, \enemyWallGuard{}s and \enemyCultist{}s.
	\item The parshmen can have different coloration of patterns and occupation.
	\item If the party is quick to follow queen's order, introduce moral dilemma by making subsequent parshmen even more helpless: make them old and barely moving, pregnant or just a kid.
\end{itemize}



\subsection{Hostile Rioters}

When you start this encounter, read the following:

\quoteBox{
	You did not see that coming. A heavy blow hits you in the back. "You are one of them, aren't you?", the rioter asks. "Where do you think you are going? This street is ours!"
}

The party is attacked by 1d6 + 2 \enemyRioter{}s. Each character who doesn't pass \skillCheck{15}{Perception} is Surprised in the first round of the fight. The enemies fight until there are only 2 of them left, at which moment they flee.



\subsection{Cult Procession}

1d4 + 4 \enemyCultist{}s (each \enemyCultistOfMoments{} of random kind) walk in procession through Kholinar streets. They seem oblivious to riots all around them and they will not attack if unprovoked. If the characters start the fight, they fight to the end. None of them carries anything valuable.

\spoilerBox{
	Similar procession is described in \bookOB{}, Chapter 61: Nightmare Made Manifest.
}




\subsection{Hungry Beggar}

When the party meets the beggar, read the following:

\quoteBox{
	You see an old and dirty man wearing rags, who begs on the side of the street. 
	“Please, I haven’t eaten for days”, he asks, reaching his only healthy hand towards you.
}

The characters can share the food with the beggar. It is a right thing to do, but they cannot expect anything in return, and food is scarce in the city...

You can repeat this encounter several times, changing the beggar NPC each time to someone new.



\subsection{Wounded Enemy}

When you introduce this encounter, read the following:

\quoteBox{
	You can see the rioters fought here not so long ago. The street is painted red with pools of blood, dead bodies lie scattered everywhere. You want to leave when you hear faint voice: "Please... help me..."
}

The party meets one badly wounded adversary (one from \enemyWallGuard{}, \enemyRioter{} and \enemyCultistOfMoments{}). They can help him if they so choose.

\reasonBox{
	To make this encounter more interesting, choose an NPC from the fraction the characters are the most hostile towards.
}



\subsection{Trouble at Home}

If the party decides to return to \npcGrandma{} during the riots, read the following:

\quoteBox{
	You hear grandma \npcGrandmaName{} even before you enter the building: \\ \\
	"Hey! Where are your manners! Just come here, you storming son of a chul! I will teach you! Hey, I am talking to you, you hear me? Come here! You will regret..."
}

Two \enemyRioter{}s and one \enemyWallGuard{} are looting the place. They ignore the old lady and the parshmen, but they will turn hostile when they notice the PCs entering. When only one of them is left standing, he will try to flee the building.



\subsection{It's a Hard Work}

The characters may decide that the riots are the best opportunity for some side job in the goods liberation and redistribution business. Depending on the value of the loot they are aiming for, consider no resistance, some competition in a form of a few other \enemyRioter{}s, hired guards protecting the place (use \enemyWallGuard{} stat block), or other adversaries from \bookWorldGuide{}.

\reasonBox{
	Try to ensure the party will not spend entire chapter just looting random places in the city. If you feel this activity is repeated too many times, let the spren urge them to focus on their goals, increase security of looted places, or, if that wasn't enough, trigger the events of \fullref{sec:enemyReborn}.
}